\documentclass[11pt]{article}

\usepackage[a4paper,margin=25mm]{geometry}
\usepackage{amsmath,amssymb,amsthm}
\usepackage{booktabs}
\usepackage{microtype}
\usepackage[hidelinks]{hyperref}

\newcommand{\Q}{\mathbf Q}
\newcommand{\F}{\mathbf F}
\newcommand{\Pone}{\mathbf P^1}
\newcommand{\Pic}{\operatorname{Pic}}
\newcommand{\Jac}{\operatorname{Jac}}
\newcommand{\gon}{\operatorname{gon}}
\newcommand{\red}{\operatorname{red}}
\newcommand{\OO}{\mathcal O}
\newcommand{\code}[1]{\texttt{#1}}

\newtheorem{theorem}{Theorem}
\newtheorem{proposition}[theorem]{Proposition}
\newtheorem{lemma}[theorem]{Lemma}

\hypersetup{
  pdftitle={A lower bound for the rational gonality of 43.132.9.a.2},
  pdfauthor={Computational proof note}
}

\title{\vspace{-2em}\textbf{A lower bound for the rational gonality of
\(X_\Delta(43)\)}\\[3pt]
\large LMFDB label \code{43.132.9.a.2}}
\author{}
\date{}

\begin{document}
\maketitle
\vspace{-3em}

\begin{abstract}
Let \(\Delta=\langle2\rangle\subset(\mathbf Z/43\mathbf Z)^\times\).
We give a short computational proof that
\(\gon_{\Q}(X_\Delta(43))\geq7\).  The proof uses the unique
degree-\(6\) pencil class modulo \(3\) and an exact description of the
image of \(J(\Q)\) in \(J(\F_3)\).  All computations are exact and the
small certificate needed to reproduce them is specified at the end.
\end{abstract}

\section{Setup and the finite-field obstruction}

Let \(X=X_\Delta(43)\), with LMFDB label
\code{43.132.9.a.2}\footnote{\url{https://beta.lmfdb.org/ModularCurve/Q/43.132.9.a.2/}},
and let \(J=\Jac(X)\).  The curve has genus \(9\), three rational cusps
\(c_1,c_2,c_3\), and one further Galois orbit of cusps whose sum is a
rational divisor \(c_{(3)}\) of degree \(3\).  These are all the cusps.
The identifications are checked on the canonical model by pulling back
the two poles of the \(j\)-map on \(X_0(43)\): the two fibres are
\[
                     c_1+c_2+c_3
             \quad\text{and}\quad c_{(3)}.
\tag{1.1}\label{eq:cusp-fibres}
\]

The integral canonical equations have the same Hilbert polynomial
\(16n-8\) before and after reduction modulo \(3\).  The reduced curve is
absolutely irreducible and has arithmetic and geometric genus \(9\).
It is therefore smooth: for an integral curve over the perfect field
\(\F_3\), the difference of these two genera is the sum of the
nonnegative singularity invariants \(\delta_P\).  We write this smooth
reduction as \(X_3\).

\begin{proposition}\label{prop:f3}
The following hold.
\begin{enumerate}
\item \(\gon_{\F_3}(X_3)=6\).
\item \(W^1_6(X_3)(\F_3)\) consists of one class.
\item Using \(6c_3\) to identify \(\Pic^6(X_3)\) with \(J(\F_3)\), this
      class is
      \[
                         T=[1,1,1977,0].
      \tag{1.2}\label{eq:T}
      \]
\end{enumerate}
\end{proposition}

\begin{proof}
The first assertion is the exact CurveArith gonality computation
\cite{DerickxTerao}.  For the other two, the numbers of places of degrees
\(1,\ldots,5\) are
\[
\begin{array}{c|rrrrr}
\text{degree}&1&2&3&4&5\\ \hline
\text{number of places}&4&3&4&24&57.
\end{array}
\tag{1.3}\label{eq:place-counts}
\]
Every degree-\(6\) class with at least two sections has an effective
representative containing a rational place: evaluate its space of
sections at any of the four rational places.  The kernel of this
evaluation map is nonzero.

It is therefore enough to enumerate effective degree-\(6\) divisors
containing a rational place.  Repeated places are included.  Encode such
a divisor by writing its total degree-one multiplicity as \(m_1+1\), and
its number of degree-\(i\) places as \(m_i\) for \(i\geq2\).  Their number
is
\[
\sum_{\substack{m_1+2m_2+\cdots+5m_5=5\\m_i\geq0}}
 \binom{4+m_1}{m_1+1}
 \prod_{i=2}^{5}\binom{N_i+m_i-1}{m_i}
 =845,
\]
where \((N_2,N_3,N_4,N_5)=(3,4,24,57)\).
Exact Riemann--Roch calculations show that four of these divisors have
at least two sections, and all four are linearly equivalent.

Magma returns
\[
\Pic(X_3)\simeq
 \mathbf Z/2\mathbf Z\oplus\mathbf Z/2\mathbf Z
 \oplus\mathbf Z/4788\mathbf Z\oplus\mathbf Z.
\]
In this basis the unique translated degree-zero class is
\eqref{eq:T}.  The calculation also verifies that the chosen origin is
the reduction of \(6c_3\), so no translation ambiguity remains.
\end{proof}

\section{The reduction of the rational Jacobian}

There is an explicit degree-\(6\) morphism
\[
             \pi:X\longrightarrow
             E:\quad y^2+y=x^3+x^2,
\tag{2.1}\label{eq:pi}
\]
where \(E\) is the elliptic curve \(43\mathrm a1\).  Its equations and
degree are checked directly on the canonical model.

Let \(C\subset J(\Q)\) be the subgroup generated by
\[
      [c_2-c_1],\qquad[c_3-c_1],\qquad[c_{(3)}-3c_3].
\tag{2.2}\label{eq:Cgens}
\]
It is finite by Manin--Drinfeld \cite{Drinfeld}.  Let \(\overline C\)
denote its reduction in \(J(\F_3)\), and set
\[
 H_3=\langle\,\pi^*E(\F_3),\,\overline C\,\rangle
      \subset J(\F_3).
\tag{2.3}\label{eq:H3}
\]
Exact divisor-class arithmetic gives
\[
 \#E(\F_3)=6,\qquad
 \#\overline C=532,\qquad
 \#H_3=3192,\qquad
 T\notin H_3.
\tag{2.4}\label{eq:finite-membership}
\]

\begin{proposition}\label{prop:saturation}
One has
\[
                       J(\Q)=\pi^*E(\Q)+C.
\tag{2.5}\label{eq:saturation}
\]
Consequently \(\red_3J(\Q)\subset H_3\).
\end{proposition}

\begin{proof}
\noindent\emph{The elliptic factor.}
Put \(A=\operatorname{im}(\pi^*:E\to J)\).  First, \(\pi^*\) is
injective.  Indeed, if \(K=\ker(\pi^*)\neq0\), write
\(\pi^*=i\circ q\), where \(q:E\to E/K\) is the quotient isogeny.
Dualizing gives \(\pi_*=\widehat q\circ\widehat i\).  After choosing a
rational cusp as base point, \(\pi-\pi(c_1)\) is the composite of the
Abel--Jacobi map with \(\pi_*\); hence, up to translation, \(\pi\)
factors through \(\widehat q\).  It follows that
\(\#K=\deg\widehat q\) divides \(\deg\pi=6\).
Such a kernel would be cyclic.  Exact elliptic-curve computations show
that \(E\) has no nontrivial rational cyclic isogeny and that
\(E(\Q)_{\rm tors}=0\).  Thus \(A\simeq E\) and \(A(\Q)\) is
torsion-free.

\smallskip
\noindent\emph{The complementary quotient.}
Let \(B=J/A\).  Magma's modular-symbols variety
\(\code{JH(43,[2])}\) is only \(\Q\)-isogenous to \(J\), which is enough
for the following argument.  Its simple factors have dimensions and
exact \(L\)-ratios
\[
\begin{array}{c|rrrr}
\text{dimension}&1&2&2&4\\ \hline
L\text{-ratio at }s=1&0&2/7&3/16&1/361.
\end{array}
\tag{2.6}\label{eq:Lratios}
\]
The unique elliptic factor is \(43\mathrm a\), so the sum of the other
three factors is \(\Q\)-isogenous to \(B\).  Their \(L\)-values at \(1\)
are nonzero.  Kato's theorem \cite[Corollary 14.3]{Kato} therefore gives
\[
                         \operatorname{rank}B(\Q)=0.
\tag{2.7}\label{eq:rankzero}
\]

Isogenous abelian varieties have the same good-prime point counts.  The
primes below are good (the only bad prime is \(43\)), and the exact
Frobenius computations give
\[
\begin{array}{c|rrr}
p&3&5&23\\ \hline
\#B(\F_p)&3192&80864&86398048604.
\end{array}
\]
Reduction on rational torsion is injective at these good primes, and
\[
             \gcd(3192,80864,86398048604)=532.
\]
Together with \eqref{eq:rankzero}, this shows
\[
                         \#B(\Q)\leq532.
\tag{2.8}\label{eq:Bupper}
\]

\smallskip
\noindent\emph{Saturation.}
On the other hand, \(\#\overline C=532\) implies \(\#C\geq532\).
Since \(C\) is torsion and \(A(\Q)\) is torsion-free,
\(C\cap A(\Q)=0\); hence \(C\) injects into \(B(\Q)\).
Equation \eqref{eq:Bupper} forces its image to be all of \(B(\Q)\).
For every \(x\in J(\Q)\), subtracting a suitable \(c\in C\) therefore
leaves an element of \(A(\Q)=\pi^*E(\Q)\).  This proves
\eqref{eq:saturation}, and \eqref{eq:H3} then gives the final assertion.
\end{proof}

\section{Conclusion}

\begin{theorem}
For \(X=X_\Delta(43)\) with LMFDB label \code{43.132.9.a.2},
\[
                              \gon_\Q(X)\geq7.
\]
\end{theorem}

\begin{proof}
Suppose first that \(X\) has a rational map to \(\Pone\) of degree at
most \(5\).  Specializing the associated line bundle to the smooth model
over \(\mathbf Z_3\), upper semicontinuity gives a line bundle of the same
degree with at least two sections on \(X_3\).  This contradicts
\(\gon_{\F_3}(X_3)=6\).

It remains to exclude degree \(6\).  Let
\(\mathcal L=g^*\OO_{\Pone}(1)\).  The relative Picard scheme is proper,
so \(\mathcal L\) extends over \(\mathbf Z_3\), and upper semicontinuity
gives
\[
           \deg\mathcal L_3=6,\qquad
           h^0(X_3,\mathcal L_3)\geq2.
\]
By Proposition \ref{prop:f3}, after translating by \(-6c_3\), its
specialization is the unique class \(T\).  But
\([\mathcal L(-6c_3)]\in J(\Q)\), so Proposition
\ref{prop:saturation} implies that its reduction lies in \(H_3\).
This contradicts \(T\notin H_3\) in
\eqref{eq:finite-membership}.  Hence no rational map of degree at most
\(6\) exists.
\end{proof}

\section*{Reproducibility}

The canonical model and maps are the LMFDB data for
\code{43.132.9.a.2} \cite{LMFDB}.  The following Magma \(2.29\)-\(4\)
scripts, run from the repository root, verify every numerical assertion
above:
\begin{quote}\small\ttfamily
gonality\_43\_132\_9\_a\_2/audit\_smooth\_reduction.m\\
gonality\_43\_132\_9\_a\_2/audit\_certificate.m\\
gonality\_43\_132\_9\_a\_2/audit\_torsion\_counts.m\\
gonality\_43\_132\_9\_a\_2/prove\_lower\_bound\_7.m
\end{quote}
The last script performs the exhaustive divisor enumeration, computes
\(H_3\), tests \(T\notin H_3\), and verifies the modular-factor,
elliptic-isogeny, and torsion assertions.

\begin{thebibliography}{9}

\bibitem{DerickxTerao}
M.~Derickx and K.~Terao,
\emph{Computing class groups and gonalities of algebraic curves over
finite fields}, arXiv:2602.17417 (2026).

\bibitem{Kato}
K.~Kato,
\emph{\(p\)-adic Hodge theory and values of zeta functions of modular
forms}, Astérisque \textbf{295} (2004), 117--290.

\bibitem{Drinfeld}
V.~G.~Drinfeld,
\emph{Two theorems on modular curves},
Functional Anal. Appl. \textbf{7} (1973), 155--156.

\bibitem{LMFDB}
The LMFDB Collaboration,
\emph{The L-functions and Modular Forms Database},
\url{https://www.lmfdb.org/}.

\bibitem{Magma}
W.~Bosma, J.~Cannon, and C.~Playoust,
\emph{The Magma algebra system I: The user language},
J. Symbolic Comput. \textbf{24} (1997), 235--265.

\end{thebibliography}

\end{document}
